\documentclass[12pt]{article}
\usepackage{amsmath,amssymb,amsfonts}
\usepackage[margin=1in]{geometry}
\begin{document}
\begin{center}
{\large\bfseries 24th Irish Mathematical Olympiad\\7 May 2011}
\end{center}
\bigskip\noindent\textbf{Paper 1}\bigskip
\begin{enumerate}
\item Suppose $abc \ne 0$. Express in terms of $a$, $b$, and $c$ the solutions $x$, $y$, $z$, $u$, $v$, $w$ of the equations
\[
x + y = a, \quad z + u = b, \quad v + w = c, \quad ay = bz, \quad ub = cv, \quad wc = ax.
\]

\item Let $ABC$ be a triangle whose side lengths are, as usual, denoted by $a = |BC|$, $b = |CA|$, $c = |AB|$. Denote by $m_a$, $m_b$, $m_c$, respectively, the lengths of the medians which connect $A$, $B$, $C$, respectively, with the centres of the corresponding opposite sides.

\begin{enumerate}
\item Prove that $2m_a < b + c$. Deduce that $m_a + m_b + m_c < a + b + c$.

\item Give an example of
\begin{enumerate}
\item[(i)] a triangle in which $m_a > \sqrt{bc}$;
\item[(ii)] a triangle in which $m_a \le \sqrt{bc}$.
\end{enumerate}
\end{enumerate}

\item The integers $a_0, a_1, a_2, a_3, \ldots$ are defined as follows:
\[
a_0 = 1, \quad a_1 = 3, \quad \text{and } a_{n+1} = a_n + a_{n-1} \text{ for all } n \ge 1.
\]
Find all integers $n \ge 1$ for which $na_{n+1} + a_n$ and $na_n + a_{n-1}$ share a common factor greater than $1$.

\item The incircle $\mathcal{C}_1$ of triangle $ABC$ touches the sides $AB$ and $AC$ at the points $D$ and $E$, respectively. The incircle $\mathcal{C}_2$ of the triangle $ADE$ touches the sides $AB$ and $AC$ at the points $P$ and $Q$, and intersects the circle $\mathcal{C}_1$ at the points $M$ and $N$. Prove that

\begin{enumerate}
\item the centre of the circle $\mathcal{C}_2$ lies on the circle $\mathcal{C}_1$.

\item the four points $M$, $N$, $P$, $Q$ in appropriate order form a rectangle if and only if twice the radius of $\mathcal{C}_1$ is three times the radius of $\mathcal{C}_2$.
\end{enumerate}

\item In the mathematical talent show called ``The $X^2$-factor'' contestants are scored by a panel of $8$ judges. Each judge awards a score of $0$ (`fail'), $X$ (`pass'), or $X^2$ (`pass with distinction'). Three of the contestants were Ann, Barbara and David. Ann was awarded the same score as Barbara by exactly $4$ of the judges. David declares that he obtained different scores to Ann from at least $4$ of the judges, and also that he obtained different scores to Barbara from at least $4$ judges.

In how many ways could scores have been allocated to David, assuming he is telling the truth?
\end{enumerate}

\newpage
\begin{center}
{\large\bfseries 24th Irish Mathematical Olympiad\\7 May 2011}
\end{center}
\bigskip\noindent\textbf{Paper 2}\bigskip
\begin{enumerate}
\setcounter{enumi}{5}
\item Prove that
\[
\frac{2}{3} + \frac{4}{5} + \cdots + \frac{2010}{2011}
\]
is not an integer.

\item In a tournament with $N$ players, $N < 10$, each player plays once against each other player scoring $1$ point for a win and $0$ points for a loss. Draws do not occur. In a particular tournament only one player ended with an odd number of points and was ranked fourth. Determine whether or not this is possible. If so, how many wins did the player have?

\item $ABCD$ is a rectangle. $E$ is a point on $AB$ between $A$ and $B$, and $F$ is a point on $AD$ between $A$ and $D$. The area of the triangle $EBC$ is $16$, the area of the triangle $EAF$ is $12$ and the area of the triangle $FDC$ is $30$. Find the area of the triangle $EFC$.

\item Suppose that $x$, $y$ and $z$ are positive numbers such that
\[
1 = 2xyz + xy + yz + zx. \tag{1}
\]
Prove that
\begin{enumerate}
\item[(i)] $\dfrac{3}{4} \le xy + yz + zx < 1$.

\item[(ii)] $xyz \le \dfrac{1}{8}$.
\end{enumerate}
Using (i) or otherwise, deduce that
\[
x + y + z \;\ge\; \frac{3}{2}, \tag{2}
\]
and derive the case of equality in (2).

\item Find with proof all solutions in nonnegative integers $a$, $b$, $c$, $d$ of the equation
\[
11^a 5^b - 3^c 2^d = 1.
\]
\end{enumerate}
\end{document}
